\section{Jerusalem, Judea, Caesarea}

For Peter's next dozen years the almost exclusive source is the Acts of the
Apostles, and its character must govern its use. Acts is a canonical
theological narrative written a generation later, whose speeches are
compositions in character after the manner of ancient historiography, and
whose selection serves a thesis---the Spirit-driven spread of the word from
Jerusalem toward Rome. Read within those limits, its portrait of Peter is
coherent and, at key points, independently echoed by Paul.

Peter opens every decisive scene of the Jerusalem church. He initiates the
replacement of Judas (Acts 1:15--26). At Pentecost, ``standing up with the
eleven,'' he preaches the first sermon---``God hath made both Lord and
Christ, this same Jesus, whom you have crucified\ldots\ Do penance, and be
baptized every one of you'' (Acts 2:14, 36, 38)---and Luke counts ``about
three thousand souls'' (2:41). At the Beautiful Gate he heals the lame
beggar: ``Silver and gold I have none; but what I have, I give thee: In the
name of Jesus Christ of Nazareth, arise, and walk'' (Acts 3:6). Arrested,
he answers the Sanhedrin with the stone-rejected oracle and the claim that
``there is no other name under heaven given to men, whereby we must be
saved'' (Acts 4:11--12). Within the community his authority is fearsome in
the Ananias and Sapphira episode (Acts 5:1--11), and popular devotion seeks
even his passing shadow (5:15). With John he is sent to Samaria to confirm
Philip's converts, where he rebukes Simon the magician's attempt to buy the
Spirit: ``Keep thy money to thyself, to perish with thee\ldots\ thy heart
is not right in the sight of God'' (Acts 8:20--21)---the seed, as it
happens, both of the word \emph{simony} and of a vast later legend cycle.

The theological center of the Acts portrait is the Cornelius affair. After
a vision at Joppa overturning the clean--unclean boundary, Peter enters a
Gentile officer's house at Caesarea: ``In very deed I perceive, that God is
not a respecter of persons'' (Acts 10:34); the Spirit falls on the
uncircumcised household while he is still speaking (10:44), and he draws
the conclusion---``Can any man forbid water, that these should not be
baptized?'' (10:47)---then defends it before critics in Jerusalem
(11:1--18). Luke has thus made Peter, not Paul, the door through which the
Gentiles first enter. The sequel is persecution: about the year 44 Herod
Agrippa~I ``killed James, the brother of John, with the sword'' and,
``seeing that it pleased the Jews,'' seized Peter (Acts 12:1--3), who was
``sleeping between two soldiers, bound with two chains'' when the angel
struck him awake (12:6--7). Freed, he leaves a message ``to James, and to
the brethren,'' and---in five words that have carried a great weight of
speculation---``going out, he went into another place'' (12:17). Luke never
says where. Peter reappears once, at the Jerusalem council of about 48--49,
recalling Cornelius and rejecting ``a yoke\ldots\ which neither our fathers
nor we have been able to bear'': ``by the grace of the Lord Jesus Christ,
we believe to be saved, in like manner as they also'' (Acts 15:7--11).
James presides and formulates; Peter testifies first. After that verse the
Acts of the Apostles never mentions him again, and the continuous narrative
of Peter's life ends---eighteen or twenty years before his death.
